% Nouveau EGA V (cote 25): archive cote 25, site page 62 (PDF page 63).
% typescript. Draft compared: none.
% AI-transcribed by Claude Opus 5.5 (claude-opus-5-5, adaptive thinking, effort high), Message Batches API, batch ega25b (msgbatch_01A7n5YuXubUPNK1zsepi8RG), 2026-11-13; not yet checked by eye.
\scanpage{25/62}
\mspage{57}

génériquement étale\var{\inserted[typed]{de degré \hand{d}}}{de degré \hand{d}}{insertion}{ega25_25_062-v0001}, et pour tout point géométrique  s de $P^1$ ,  $X_s$
est étale sur le corps \hand{$-$}\unsure{alg} clos \var{\inserted{k'=}}{k'=}{insertion}{ega25_25_062-v0002} k(s), ou  \var{\xout{\illegible{3}}}{}{typewriter strike-out}{ega25_25_062-v0003} est
\var{\inserted{k'-}}{k'-}{insertion}{ega25_25_062-v0004}isomorphe à la somme de \var{\over{\unsure{n}}{d}}{d}{writing over}{ega25_25_062-v0005}-2  schémas \var{\inserted{Spec(k')}}{Spec(k')}{insertion}{ega25_25_062-v0006} \var{\xout{k(s)}}{}{typewriter strike-out}{ega25_25_062-v0007}  et \var{\xout{un}}{}{typewriter strike-out}{ega25_25_062-v0008}  du schéma  \var{\xout{k}}{}{typewriter strike-out}{ega25_25_062-v0009}
$I_{k'} = \mathrm{Spec}\, k'[t]/(t^2)$ .

En langage de Papa, au  dessus de s  il y a au plus un point de
ramification, et la ramification en c e point est ``quadratique'' .

11.  \uline{Grassmaniennes}. \var{\deleted{Ratrapages.}}{}{deletion}{ega25_25_062-v0010}

Comme nous avions couramment à utiliser des sous-variétés linéai-%
res de P , pas seulement de dimension relative 0 ou r-1, il est clair
qu'on aura besoin de quelques notions sur les grassmaniennes  et\var{\inserted{de}}{de}{insertion}{ega25_25_062-v0011}quelques
sorties de nature ``géométrie élémentaire'' sur les constructions concer-%
nant les variétés linéaires, qui  devraient venir tout au début du par.
De plus, on aura en pratique à prendre parfois des sections linéaires
quelcohques, pas seulement des sections hyperplanes, et il y a lieu de
revoir dans cet esprit élargi tous les n°s antérieurs.

Soit  $\underline{E}$  un Module qu coh  sur le préschéma $\unsure{S|X}$\var{\inserted[typed]{et n un entier \hand{$>$} 0.}}{et n un entier \hand{$>$} 0.}{insertion}{ega25_25_062-v0012}\hand{,} Considérons le
foncteur \qquad $(Sch)_{/S}^{\,o} \longrightarrow (Ens)$  défini par
\[ \underline{Grass}_n(\underline{E})(S') = \text{ ensemble des Modules quotients loc libres} \]
\[ \text{de rang  n  de } \underline{E}_{S'} \ . \]
Ce foncteur est représentable, et le préschéma sur S qui le réprésente
sera donc noté également $\underline{Grass}_n(\underline{E})$.  Pour prouver la r eprésentabilité,
considérer l'homomorphisme de foncteurs naturels
\[ \underline{Grass}_n(\underline{E}) \longrightarrow \underline{Grass}_1( \overset{\hand{n}}{\textstyle\bigwedge} \underline{E} ) = \underline{P} ( \overset{\hand{n}}{\textstyle\bigwedge} \underline{E}) \ , \]
défini en associant à  tout  quotient loc libre  de rang n  $\underline{G}$  de $\underline{E}_{S'}$
le Module  loc libre de rang 1 \quad $\overset{n}{\textstyle\bigwedge} \underline{G}$ , considéré comme un quotient
de  $\overset{n}{\textstyle\bigwedge} \underline{E}_{S'}$ \hand{\unsure{?}}  On prouve comme dans  Séminaire Cartan que  ce morphisme
est  ``représentable par \var{\xout{\illegible{1}}}{}{typewriter strike-out}{ega25_25_062-v0013} immersion fermées'' , de sorte que $\underline{Grass}_n(\underline{E})$

\mspage{62}
